% Einheit 2 von 4 – Proportionale Zuordnungen und Dreisatz (bank/zuordnungen/e2.jsonl, zusammenbau v0.1)
%% TODO Zweigzeile Teil 1: was hier gelernt wird (Satz in Schülersprache aus dem Einheitstitel) – Einheit 2
\einheitenkopf[e2]{Einheit 2 von 4 · Proportionale Zuordnungen und Dreisatz}
\zweigzeile{ab Kl. 6 · P10}
\setcounter{aufgabe}{11}

%% TODO Titel: Ich-kann-Satz fehlt in der Bank (Platzhalter: Kettenname) – Nr. 12
%% TODO Anweisung: ein Satz über der ersten Teilaufgabe fehlt in der Bank – Nr. 12
\begin{aufgabe}{Was ist hier eine Portion?}
\begin{teile}
\teil Äpfel: 3 kg kosten 4,50 €. Was ist hier eine Portion?
\end{teile}
\end{aufgabe}

%% TODO Titel: Ich-kann-Satz fehlt in der Bank (Platzhalter: Kettenname) – Nr. 13
%% TODO Anweisung: ein Satz über der ersten Teilaufgabe fehlt in der Bank – Nr. 13
\begin{aufgabe}{Minitabelle anlegen}
\begin{teile}
\teil 5 kg Kartoffeln kosten 6 €. Was kosten 8 kg? Rechne noch nicht: Welche Mengen stehen in der Minitabelle links untereinander?
\end{teile}
\end{aufgabe}

%% TODO Titel: Ich-kann-Satz fehlt in der Bank (Platzhalter: Kettenname) – Nr. 14
%% TODO Anweisung: ein Satz über der ersten Teilaufgabe fehlt in der Bank – Nr. 14
\begin{aufgabe}{Pfeile beidseitig}
\begin{teile}
\teil 8 Hefte kosten 10 €. Im Dreisatz steht links am Pfeil „: 8“. Was steht rechts am Pfeil?
\end{teile}
\end{aufgabe}

%% TODO Titel: Ich-kann-Satz fehlt in der Bank (Platzhalter: Kettenname) – Nr. 15
%% TODO Anweisung: ein Satz über der ersten Teilaufgabe fehlt in der Bank – Nr. 15
\begin{aufgabe}{Dreisatz proportional}
\swa{Führe die Tabelle weiter.}{\wertetabelle[2,4,,]{Anzahl Stück}{Preis in €}{1,2,3,4}}
%% TODO Darstellung für schwach fehlt in der Bank (zuordnungen-e2-k4-s1-v1; Abschnitt „Für schwache Schüler“)
\swz{2 kg Birnen kosten 3 €. Was kosten 4 kg?}{}
%% TODO Darstellung für schwach fehlt in der Bank (zuordnungen-e2-k4-s1-v2; Abschnitt „Für schwache Schüler“)
\swz{6 Eier kosten 2 €. Was kosten 3 Eier?}{}
%% TODO Darstellung für schwach fehlt in der Bank (zuordnungen-e2-k4-s1-v3; Abschnitt „Für schwache Schüler“)
\swz{Ein Radfahrer schafft in 10 Minuten 3 km. Wie weit kommt er in 20 Minuten?}{}
%% TODO Darstellung für schwach fehlt in der Bank (zuordnungen-e2-k4-s1-v4; Abschnitt „Für schwache Schüler“)
\swz{8 Brötchen kosten 4 €. Was kosten 4 Brötchen?}{}
%% TODO Darstellung für schwach fehlt in der Bank (zuordnungen-e2-k4-s1-v5; Abschnitt „Für schwache Schüler“)
\swz{5 l Farbe reichen für 15 m² Wand. Wie viel Wand schafft man mit 10 l?}{}
\swfrage{Was bleibt gleich, was ändert sich?}
%% TODO Darstellung für schwach fehlt in der Bank (zuordnungen-e2-k4-s2-v1; Abschnitt „Für schwache Schüler“)
\swz{5 Joghurts kosten 3 €. Was kosten 15 Joghurts?}{}
%% TODO Darstellung für schwach fehlt in der Bank (zuordnungen-e2-k4-s3-v1; Abschnitt „Für schwache Schüler“)
\swz{30 Hefte kosten 45 €. Was kosten 10 Hefte?}{}
%% TODO Darstellung für schwach fehlt in der Bank (zuordnungen-e2-k4-s4-v1; Abschnitt „Für schwache Schüler“)
\swz{3 Kugeln Eis kosten 6 €. Was kosten 7 Kugeln?}{}
%% TODO Darstellung für schwach fehlt in der Bank (zuordnungen-e2-k4-s5-v1; Abschnitt „Für schwache Schüler“)
\swz{5 Brezeln kosten 3,50 €. Was kosten 9 Brezeln?}{}
\begin{teile}
\teil Packung A: 500 g Nudeln für 2,00 €. Packung B: 750 g für 2,70 €. Welche ist günstiger? \\ \kreuz{Packung A} \\ \kreuz{Packung B}
\end{teile}
\swz[3]{Ein Läufer braucht für 2 km 13 Minuten. Wie lange braucht er bei gleichem Tempo für 5 km? Rechne ohne Taschenrechner.}{}
\end{aufgabe}

%% TODO Titel: Ich-kann-Satz fehlt in der Bank (Platzhalter: Kettenname) – Nr. 16
%% TODO Anweisung: ein Satz über der ersten Teilaufgabe fehlt in der Bank – Nr. 16
\begin{aufgabe}{fester Faktor angeben (Preis je Einheit), Gleichung y = k · x}
%% TODO Darstellung für schwach fehlt in der Bank (zuordnungen-e2-k5-s1-v1; Abschnitt „Für schwache Schüler“)
\swz{7 kg Äpfel kosten 17,50 €. Gib den Preis je kg an und schreibe die Gleichung (x: Menge in kg, y: Preis in €).}{}
\end{aufgabe}

%% TODO Titel: Ich-kann-Satz fehlt in der Bank (Platzhalter: Kettenname) – Nr. 17
%% TODO Anweisung: ein Satz über der ersten Teilaufgabe fehlt in der Bank – Nr. 17
\begin{aufgabe}{Graph zeichnen (Ursprungsgerade) und Werte ablesen}
\swa{1 kg Kirschen kostet 3 €. Zeichne den Graphen Menge in kg → Preis in € und lies ab, was 2,5 kg kosten.}{\begin{ksys}[xmin=0,xmax=5,ymin=0,ymax=15,xlabel=Menge in kg,ylabel=Preis in €]\end{ksys}}
\end{aufgabe}

%% TODO Titel: Ich-kann-Satz fehlt in der Bank (Platzhalter: Kettenname) – Nr. 18
%% TODO Anweisung: ein Satz über der ersten Teilaufgabe fehlt in der Bank – Nr. 18
\begin{aufgabe}{Verhältnisgleichung aufstellen (Vorrat)}
%% TODO Darstellung für schwach fehlt in der Bank (zuordnungen-e2-k7-s1-v1; Abschnitt „Für schwache Schüler“)
\swz{$\text{Für 4 Portionen Pfannkuchenteig braucht man 300 g Mehl. Wie viel Mehl braucht man für 10 Portionen? Stelle eine Verhältnisgleichung auf und löse sie (x: Mehl in g).}$}{}
\end{aufgabe}

%% TODO Titel: Ich-kann-Satz fehlt in der Bank (Platzhalter: Kettenname) – Nr. 19
%% TODO Anweisung: ein Satz über der ersten Teilaufgabe fehlt in der Bank – Nr. 19
\begin{aufgabe}{Dreisatz proportional – Fehler finden, Begründen, Darstellungswechsel, Anwendung}
\swz[3]{3 kg Kartoffeln kosten 7,20 €. Tim rechnet für 1 kg: 7,20 € − 3 = 4,20 €. Finde den Fehler und rechne richtig.}{}
\swfrage{Warum rechnet man im Dreisatz rechts auch „: 4“, wenn man links „: 4“ rechnet? Begründe.}
\swz{Der Graph zeigt Menge in kg → Preis in €. Lies den Preis für 1 kg ab und schreibe die Gleichung (x: Menge in kg, y: Preis in €).}{\begin{ksys}[xmin=0,xmax=4,ymin=0,ymax=6,ablesen,xlabel=Menge in kg,ylabel=Preis in €]\gerade{1.5}{0}{}\end{ksys}}
\swz[3]{Ein Rezept für 4 Personen braucht 250 g Nudeln. Zum Geburtstag kommen 10 Personen. Reicht eine Packung mit 500 g?}{}
\end{aufgabe}

%% TODO Merkkasten Einheit 2: 6 Zeilen, Vorgabe höchstens fünf (3.1)
\uebersichtskasten{Proportional: Pro Portion kommt immer das Gleiche dazu – doppelte Menge, doppelter Preis; 0 Stück kosten 0 €. \\ Dreisatz: erst auf eine Portion runterrechnen, dann hochrechnen – links und rechts dasselbe rechnen. \\ 7 Stück kosten 5,60 € \quad | : 7 \\ 1 Stück kostet 0,80 € \quad | · 11 \\ 11 Stück kosten 8,80 € \\ Fester Faktor (Preis je Portion): 5,60 : 7 = 0,80 → Preis = 0,80 · Stückzahl}
